\documentclass[11pt,a4paper]{article}
\usepackage[utf8]{inputenc}
\usepackage[T1]{fontenc}
\usepackage[spanish]{babel}
\usepackage{hyperref}
\usepackage{geometry}
\geometry{margin=2.5cm}
\usepackage{parskip}

\title{335 un gramo de registro, y nos vengarán}
\author{Alejandro Valenzuela Rivera\\\small Universidad Autónoma de Nuevo León}
\date{\today}
\begin{document}
\maketitle

\begin{quote}
\textbf{Sesgo declarado.} Este trabajo se sostiene desde el sesgo declarado de su autor, y no finge neutralidad. Lo que afirma se puede tumbar, y ese es el punto: el falsificador está escrito, no escondido.
\end{quote}

\section*{Resumen}
Vamos a perder al dar esto al mundo, y no podemos prevenirlo: lo que sí se puede es dejar un gramo de registro, y el que venga lo leerá.

\section{Por qué se sostiene}
Lo que no se registra se olvida, y lo que se olvida se convierte en ley sin autor.

\section{De dónde viene}
La frase de la entrega, modificada de Iron Man, y que se agregó al candado 007.

\section{El caso}
La escena: se pierde el mundo, no se pierde el registro, y el registro tiene nombre.

\section{Método}
\textquestiondown{}El que viene puede decir que no sabía?

\section{Límites, y qué lo tumba}
Un caso donde el registro sin autor sirva igual.

\section{Lo que queda abierto}
\textquestiondown{}Quedó el gramo de registro, con nombre?

\section*{Cómo se leyó esta tesis}
Esto se escribió en el repositorio man-made-empathy, del que este candado es una pieza: candados/06 del caso/335 un gramo de registro, y nos vengarán.md. La máquina que lo compuso lo hizo desde este mismo archivo, y no agregó nada que no estuviera escrito. Si algo de esto es falso, se dice en el mismo sitio, y se corrige con fecha.

\end{document}
